\documentclass{article}
\input{~/studies/preamble}

\title{Expectation-Maximization Algorithm}
\author{Hyoung Joon, Kim}
\date{\today}

\begin{document}
\maketitle
\tableofcontents
\newpage

\section{Introduction}
    We discuss about \textbf{Expectation-Maximization algorithm} (EM algorithm) in detail.
    EM algorithm is used to find the MLE or MAP of a given likelihood function with a latent variable.
    We derive and explain EM algorithm only for discrete latent variables, but EM for continuous variables can be derived in the same way,
    exchanging summation with integration.

\section{Main Idea}
    \subsection{Intuition}
        Let \textbf{X} be a dataset in which the $n$-th row represents $\textbf{x}_n^{\top}$,
        and let \textbf{Z} be lantent variables in which the $n$-th row represents $\textbf{z}_n^{\top}$.
        Representing $\bm\theta$ as the model parameter, the log likelihood function is given by
        $\log{p(\textbf{X}|\bm\theta)}=\log{\left\{ \sum_{\textbf{Z}}p(\textbf{X},\textbf{Z}|\bm\theta) \right\}}$.
        
        Maximizing the log likelihood is straightforward if all variables are observed. However, the latent variables are
        hardly known. The information we can get about latent variable is the posterior distribution $p(\textbf{Z}|\textbf{X},\bm\theta)$.
        So, instead of maximizing the whole log likelihood, we maximize the expectation of log-likelihood under the posterior distribution of 
        latent variables. Calculating the expectation is the \textbf{E step} of the EM algorithm.
        The subsequent step, the \textbf{M step} of the EM algorithm, then maximizes the expectation calculated at the E step.
        These pair of steps give rise to the revised estimate of parameters $\bm\theta$.
        
    \subsection{Algorithm}
        \begin{algorithm}[H]
            \DontPrintSemicolon
            \SetArgSty{textnormal}
            \SetKwInOut{Input}{Input}
            \SetKwInOut{Output}{Output}
            
            \caption{EM algorithm}

            \Input{joint distribution $p(\textbf{X},\textbf{Z}|\bm\theta)$ \newline
            initial parameters $\bm\theta^{(\text{old})}$ \newline
            dataset \textbf{X}}
            \Output{Final parameters $\bm\theta$}

            \BlankLine
            \hrule

            \Repeat{convergence}
            {
                \tcp{E step}
                $\mathcal{Q}(\bm\theta,\bm\theta^{(\text{old})}) \leftarrow \sum_{\textbf{Z}}p(\textbf{Z}|\textbf{X},\bm\theta^{(\text{old})})\log{p(\textbf{X},\textbf{Z}|\bm\theta)}$\;
                \tcp{M step}
                $\bm\theta^{(\text{new})} \leftarrow \arg\max_{\bm\theta}{\mathcal{Q}(\bm\theta,\bm\theta^{(\text{old})})}$\;
                $\bm\theta^{(\text{old})} \leftarrow \bm\theta^{(\text{new})}$\;
                $\mathcal{L} \leftarrow \log{p(\textbf{X}|\bm\theta)}$\;
            }
            \Return{$\bm\theta^{(\text{new})}$}
        \end{algorithm}

        If we want to use the prior distribution $p(\bm\theta)$ of parameter $\bm\theta$, we slightly modify the algorithm;
        we change M step so that it maximizes $\mathcal{Q}(\bm\theta,\bm\theta^{(\text{old})}) + \log{p(\bm\theta)}$.

    \subsection{Evidence Lower Bound}
        EM algorithm can be seen at another perspective, called \textbf{Evidence lower bound} (ELBO). The ELBO plays important role in
        modern deep generative models and \emph{variational inferences}.

        \subsubsection{Decomposition of Likelihood}
            Let $q(\textbf{Z})$ be a distribution over the latent variables. Then, the following decomposition holds for log likelihoods;
            \begin{thm}{Decomposition of log-likelihood Function}{decompLogLike}
                Let $p(\textbf{X}|\bm\theta)$ be a likelihood of dataset \textbf{X}. Let $q(\textbf{Z})$ be a distribution over the latent variables.
                Then, the following holds;
                \[
                \log{p(\textbf{X}|\bm\theta)}=\mathcal{L}(q,\bm\theta) + \text{KL}(q\|p)
                \]
                where
                \[
                \mathcal{L}(q,\bm\theta)=\sum_{\textbf{Z}}q(\textbf{Z})\log{\left\{ \frac{p(\textbf{X},\textbf{Z}|\bm\theta)}{q(\textbf{Z})} \right\}}
                \text{ and }
                \text{KL}(q\|p)=-\sum_{\textbf{Z}}q(\textbf{Z})\log{\left\{ \frac{p(\textbf{Z}|\textbf{X},\bm\theta)}{q(\textbf{Z})} \right\}}.
                \]
            \end{thm}

            Notice that the KL-divergence is always non-negative, i.e., $\text{KL}(q\|p) \geq 0$. Using this, we can derive the following useful result:
            $\mathcal{L}(q,\bm\theta)$ provides the lower bound for log-likelihood such that $\mathcal{L}(q,\bm\theta) \leq \log{p(\textbf{X}|\bm\theta)}$.

        \subsubsection{EM Revisited}
            We can derive EM algorithm using Theorem \ref{thm:decompLogLike}. We first maximize the lower bound $\mathcal{L}(q,\bm\theta^{(\text{old})})$ about $q(\textbf{Z})$,
            fixing the parameter as $\theta^{(\text{old})}$. Since the parameter is fixed, the log-likelihood stays fixed and does not fluctuate as $q(\textbf{Z})$ changes.
            So, to maximize $\mathcal{L}(q,\bm\theta^{(\text{old})})$, we minimize the KL-divergence, which can be accomplished when $q(\textbf{Z})=p(\textbf{Z}|\textbf{X},\bm\theta^{(\text{old})})$.
            We now maximize lower bound about $\bm\theta$ with fixed $q(\textbf{Z})$ this time. Since lower bound increases and the KL-divergence is non-negative since
            $q(\textbf{Z})$ is fixed while $\bm\theta$ is modified, this will increase the log-likelihood.
            Substituting $q(\textbf{Z})=p(\textbf{Z}|\textbf{X},\bm\theta^{(\text{old})})$ into the equation about lower bound of
            Theorem \ref{thm:decompLogLike}, we get
            \begin{align*}
            \mathcal{L}(q,\bm\theta) &= \sum_{\textbf{Z}}p(\textbf{Z}|\textbf{X},\bm\theta^{(\text{old})})\log{p(\textbf{X},\textbf{Z}|\bm\theta)}
            -\sum_{\textbf{Z}}p(\textbf{Z}|\textbf{X},\bm\theta^{(\text{old})})\log{p(\textbf{Z}|\textbf{X}, \bm\theta^{(\text{old})})}
            \\ &=\mathcal{Q}(\bm\theta,\bm\theta^{(\text{old})}) + \text{const}
            \end{align*}
            where the constant is independent of $\bm\theta$. Then, the E step corresponds to maximizing $\mathcal{L}(q,\bm\theta)$ about $q(\textbf{Z})$
            and the M step corresponds to maximizing $\mathcal{L}(q,\bm\theta)$ about $\bm\theta$.

\section{Extensions of EM Algorithm}
    \subsection{Generalized EM Algorithm}
        When the model gets complex, the E step or M step, or both, can be intractable. So, instead of directly maximizing $\mathcal{L}(q,\bm\theta)$ with respect to $\bm\theta$,
        we try to make the lower bound increase by modifying $\bm\theta$. For example, we can use gradient based algorithms in M step.

    \subsection{Sequential EM Algorithm}
        \emph{To be continued...}
\end{document}